\section{总结与延伸}

本节把整期压缩成五个判断。第一，Rokid 的主线不是单一产品，而是持续寻找下一代个人计算入口。第二，AI 音箱没有成为平台，是因为输入输出和场景都太窄；智能眼镜的入口机会来自 always-on 和显示能力。第三，2019 年从 AI 转 AR 是战略延续，不是突然改方向。第四，中美智能眼镜定义不同：中国近视人群和全天佩戴场景更支持 AI+AR 逻辑，美国 sunglasses/社交逻辑更强。第五，硬件创业的真正难点是产品、供应链、产能、生态和巨头竞争同时过线。

\begin{importantbox}{把 EP104 放进张小珺 AI/互联网队列}
EP104 与 EP106/EP109 的物理世界 AI 线索不同，它讲的是“随身智能”入口。若具身智能是 AI 进入外部物理世界，智能眼镜则是 AI 贴近个人感知和信息流的入口。
\end{importantbox}

\subsection{概念压缩：随身智能的三层门槛}

本节把全文再压缩成一个判断模型。智能眼镜要成为随身智能入口，至少要过三层门槛。第一层是佩戴门槛：重量、镜片、度数、外观、续航和眩晕感必须足够日常。第二层是交互门槛：它必须在短任务上明显少于手机路径成本，并且显示/语音/触控组合足够自然。第三层是生态门槛：用户时长、开发者、模型伙伴、内容场景和渠道复购要能滚动起来。

\noindent\begin{tabularx}{\textwidth}{p{0.20\textwidth}p{0.34\textwidth}X}
\toprule
门槛 & 核心问题 & EP104 给出的答案 \\
\midrule
佩戴门槛 & 用户愿不愿意一直戴 & 先从有框眼镜人群开始，不强行说服抗拒者。 \\
交互门槛 & 是否比手机更直达 & always-on、显示、语音和视线方向降低路径成本。 \\
生态门槛 & 是否能从产品走向平台 & 两小时使用、开发者、模型合作和渠道复进货共同验证。 \\
\bottomrule
\end{tabularx}

\subsection{与前后访谈的连接}

本节把 EP104 放回张小珺这一批 AI/互联网访谈。EP106 和 EP109 关注具身智能，讨论机器人如何通过数据、仿真和模型进入物理世界；EP104 则关注随身智能，讨论 AI 如何贴近个人感知和信息流。二者都不是单纯“更强模型”的问题，而是模型能力必须找到合适硬件、交互方式和数据闭环。

如果把 AI 硬件化看成一张地图，机器人是“行动端”：它替人移动、抓取、操作和服务；智能眼镜是“感知端”：它替人看、听、查询、翻译、提醒和调度注意力。二者都需要长期供应链和场景验证，但风险结构不同。机器人难在行动可靠性和物理世界成本，眼镜难在佩戴舒适、入口效率和生态密度。

\subsection{关键 takeaways}

\begin{enumerate}
\item 智能眼镜的核心不是拍照或炫技，而是 AI 的随身入口效率。
\item 显示能力会改变 AI 的交互形态，让翻译、提示和信息消费从串行语音变成可视化流程。
\item 创业公司面对巨头的窗口期有限，必须在大厂看上之前形成体验和生态壁垒。
\item 硬件爆火会反向提高质量和产能压力，不能把流量当作 PMF 本身。
\item Rokid 的案例说明，硬件创业需要长期路线、供应链能力、资本耐心和持续找问题的产品文化。
\end{enumerate}

\subsection{拓展阅读}

\begin{itemize}
\item 对 AI 如何进入物理世界感兴趣，可对照 EP106 王鹤、EP109 谢晨和 EP121 谭捷机器人访谈。
\item 对硬科技创业和供应链感兴趣，可对照 EP111 李一帆激光雷达创业史。
\item 对 AI 应用与产品入口感兴趣，可对照 EP123 ONE2X 与 EP130 AI 产品方法论访谈。
\end{itemize}

\end{document}
